\documentclass[../../main.tex]{subfiles}
\begin{document}

% SOURCE p=223; printed=208; object=ch06-notes-title; region=bottom
\section{注记与参考文献}\label{section:ch06-notes-references}\label{section:ch06-linfty-dual}

% SOURCE p=223,224; printed=208-209; object=ch06-notes-lp-history; region=bottom-top
Lebesgue 积分问世不久，人们便认识到 $L^2([a,b])$ 的重要性，因为它与 Fourier 级数及其他正交展开密切相关。Lebesgue 理论早期的一项重大成就，是 Fischer\cite{folland-bib-44}与 F.~Riesz\cite{folland-bib-114}于1907年发现 $L^2([a,b])$ 与 $\ell^2$ 同构；等价地，$L^2([a,b])$ 是完备的。F.~Riesz\cite{folland-bib-117}首先研究了 $1<p<\infty$ 时的 $L^p([a,b])$，证明了\cref{section:ch06-lp-basic,section:ch06-lp-dual}中的全部主要结果，还证明了 $L^p$ 闭单位球的弱序列紧性。$(L^1)^*=L^\infty$ 最早由 Steinhaus\cite{folland-bib-143}证明。
% SUPPLEMENT: id=v14-ch06-notes-l2-isomorphism; parent=ch06-notes-lp-history; kind=proof+motivation+hypotheses; origin=independent-derivation
\begin{remark}
\textbf{同构与完备性的含义及证明（重构新增）。}这里 $[a,b]$ 是 $a<b$ 的有限实区间，测度为 Lebesgue 测度，同构可以取为保持内积的线性满射。由\cref{theorem:ch06-lp-completeness}及内积 $\langle f,g\rangle=\int f\overline g$，$L^2([a,b])$ 是 Hilbert 空间；\cref{exercise:ch05-lebesgue-l2-separable}给出它的可分性。取可数个互不相交、长度为正的子区间，其归一化示性函数构成无限正交组，所以此空间无限维。可分 Hilbert 空间的正交基定理给出可数无限标准正交基 $(e_n)$，映射
\begin{align*}
U(f)=(\langle f,e_n\rangle)_{n\geqslant1}
\end{align*}
由 Parseval 等式保持范数；对每个 $(c_n)\in\ell^2$，部分和 $\sum_{n=1}^N c_ne_n$ 的尾差范数平方为相应的 $\sum|c_n|^2$，故由完备性收敛于一个 $f$，且逐坐标取内积得 $U(f)=(c_n)$。反过来，若已有到完备空间 $\ell^2$ 的等距满射，则将任意 Cauchy 列送到 $\ell^2$，取极限后以逆映射拉回，便得 $L^2$ 的完备性。因此原文的“等价”使用了此具体空间的可分性与无限维性，并非声称任意完备赋范空间都同构于 $\ell^2$。非退化条件 $a<b$ 排除零空间；指数二使积分范数来自内积，正交展开的系数于是精确记录范数。
\end{remark}
% END SUPPLEMENT: id=v14-ch06-notes-l2-isomorphism
% SUPPLEMENT: id=v14-ch06-notes-weak-sequences; parent=ch06-notes-lp-history; kind=proof+motivation+hypotheses; origin=independent-derivation
\begin{proposition}[有限区间上单位球的弱序列紧性（重构新增）]\label{proposition:ch06-unit-ball-weak-sequences}
若 $1<p<\infty$，则 $L^p([a,b])$ 闭单位球中的每个序列都有弱收敛的子列，其极限仍在此闭单位球中。
\end{proposition}
\begin{proof}
令 $q=p/(p-1)$，故 $1<q<\infty$。由\cref{exercise:ch06-lebesgue-separability}（将区间上的函数以零延拓至实线，再限制可数稠密集），$L^q([a,b])$ 有可数稠密集 $(g_k)$。若 $\|f_n\|_p\leqslant1$，则对每个 $k$，序列 $\int f_ng_k$ 由 Hölder 不等式有界。逐次取收敛子列，再取对角子列 $(f_{n_j})$，使这些积分对每个 $k$ 都收敛。

对任意 $g\in L^q$ 及任意 $j,\ell$，
\begin{align*}
\left|\int(f_{n_j}-f_{n_\ell})g\right|
\leqslant2\|g-g_k\|_q+
\left|\int(f_{n_j}-f_{n_\ell})g_k\right|.
\end{align*}
先取 $k$ 使第一项任意小，再取 $j,\ell$ 足够大，证明 $\int f_{n_j}g$ 是 Cauchy 列。故 $\Lambda(g)=\lim_j\int f_{n_j}g$ 对所有 $g\in L^q$ 存在，逐点取极限可知它线性，且 $|\Lambda(g)|\leqslant\|g\|_q$。由\cref{theorem:ch06-lp-duality}应用于 $L^q$，有 $f\in L^p$、$\|f\|_p\leqslant1$，满足 $\Lambda(g)=\int fg$。再由同一定理应用于 $L^p$，每个连续线性泛函都由某个 $g\in L^q$ 表示，于是 $f_{n_j}\rightharpoonup f$。
\end{proof}
\begin{remark}
\textbf{动机与逐条件说明（重构新增）。}弱收敛只检测所有线性积分配对；先在可数稠密测试集中作对角选择，再用统一范数界覆盖全部测试函数。$p>1$ 使 $q<\infty$，从而这里的可分性工具适用；$p<\infty$ 使 $(L^p)^*=L^q$ 的表示适用，排除了 $L^\infty$ 的额外泛函。单位球给出对子列和测试误差均统一的界，并使极限仍有范数至多一。区间配 Lebesgue 测度保证所用可分性和对偶定理的假设；一般空间需要另行核对这些条件。本证明直接得到序列结论，没有把弱紧性自动等同于弱序列紧性。
\end{remark}
% END SUPPLEMENT: id=v14-ch06-notes-weak-sequences

% SOURCE p=224; printed=209; object=ch06-notes-reciprocal-exponents; region=top
从某些方面看，没有把 $L^p$ 空间命名为 $L^{1/p}$ 空间，颇为可惜：正如共轭关系 $p^{-1}+q^{-1}=1$ 及\cref{section:ch06-lp-interpolation}的结果所显示的，不同 $L^p$ 空间之间的关系，通常涉及关于指数倒数的线性方程。

% SOURCE p=224; printed=209; object=ch06-notes-lieb-loss; region=upper
关于 $L^p$ 空间的更深层次性质及其在分析其他领域中的应用，可参见 Lieb 与 Loss\cite{folland-bib-93}。

% SOURCE p=224; printed=209; object=ch06-notes-basic-inequalities-history; region=upper
\textbf{关于\cref{section:ch06-lp-basic}。}Hölder 不等式在 $p=2$ 时通常与 Cauchy、Buniakovsky 和 Schwarz 的名字联系在一起：Cauchy 证明了有限和的情形，Buniakovsky 与 Schwarz 则各自独立证明了积分情形。一般 $p$ 的情形由 Hölder 与 Rogers 独立发现。Minkowski 最初的不等式是针对有限和而言的。有关文献见 Hardy、Littlewood 与 Pólya\cite{folland-bib-65}。Rubel\cite{folland-bib-122}给出了一个利用复函数论的简洁 Hölder 不等式证明。

% SOURCE p=224; printed=209; object=ch06-notes-lp-inclusions-history; region=middle
Alvarez\cite{folland-bib-5}研究了\cref{exercise:ch06-lp-sum-banach}所定义的 $L^p+L^q$ 空间之间的关系。关于\cref{exercise:ch06-lp-inclusion-criterion}以及保证 $L^p\subset L^q$ 的其他条件，可参见 Romero\cite{folland-bib-120}；关于更一般的关系 $L^p(\mu)\subset L^q(\nu)$，可参见 Miamee\cite{folland-bib-100}。

\subfile{section-06-ba.tex}

% SOURCE p=224; printed=209; object=ch06-notes-kernel-history; region=lower
\textbf{关于\cref{section:ch06-useful-inequalities}。}\cref{theorem:ch06-kernel-bound}推广了 Schur\cite{folland-bib-129}在 $p=2$ 时的结果，以及 W.~H.~Young\cite{folland-bib-164}在 $K(x,y)=k(x-y)$ 时的结果（参见\cref{section:ch08-convolution}）。\cref{theorem:ch06-homogeneous-kernel}本质上也归功于 Schur\cite{folland-bib-129}。如果本节的不等式尚未满足读者的兴趣，Hardy、Littlewood 与 Pólya\cite{folland-bib-65}的著作还提供了丰富内容。

% SOURCE p=224,225; printed=209-210; object=ch06-notes-weak-history; region=bottom-top
\textbf{关于\cref{section:ch06-distribution-weak-lp}。}弱 $L^p$ 空间最初隐含地出现在弱型估计中，这类估计的例子可以追溯到20世纪20年代；也可参见下面关于插值的注记。递减重排（\cref{exercise:ch06-decreasing-rearrangement}）由 Hardy 与 Littlewood\cite{folland-bib-064}引入。他们以板球运动的平均成绩，对其主要重排定理给出了一个有趣的动机说明。

\subfile{section-06-interpolation.tex}
\end{document}
